\begin{table}[H]
\centering
\small
\caption{Sectoral Output Granger Tests: Falsification of Austrian ABCT ($H_2$) (Chile, 1960--1980)}
\label{tab:monthly_h2_falsification}
\begin{tabular}{p{1.2cm}p{3.4cm}p{2.4cm}p{3.4cm}p{2.4cm}}
\toprule
& \multicolumn{2}{c}{\textbf{Consumer Manufacturing: $g_{M0} \to g_{\text{Manuf}}$}} & \multicolumn{2}{c}{\textbf{Capital-Intensive Mining: $g_{M0} \to g_{\text{Mining}}$}} \\
\textbf{Lag} & \textbf{$F$-Statistic} & \textbf{$p$-value} & \textbf{$F$-Statistic} & \textbf{$p$-value} \\
\midrule
1 & 0.663 & 0.4161 & 0.129 & 0.7195 (Insignificant) \\
2 & 3.636$^{**}$ & 0.0278 & 0.042 & 0.9590 (Insignificant) \\
3 & 7.149$^{***}$ & 0.0001 & 0.400 & 0.7528 (Insignificant) \\
4 & 5.256$^{***}$ & 0.0004 & 1.052 & 0.3811 (Insignificant) \\
\bottomrule
\multicolumn{5}{p{13.8cm}}{\footnotesize \textit{Notes:} Monthly physical production indices from Central Bank of Chile ($N=251$). Money growth drove consumer goods manufacturing reactivation ($F=7.15, p = 0.0001$ at lag 3) while having zero causal effect on capital-intensive mining ($p > 0.38$), directly refuting the Austrian roundabout capital malinvestment mechanism ($H_2$). *** $p < 0.01$, ** $p < 0.05$.}
\end{tabular}
\end{table}
